
Uncertainty-based hallucination detection methods report strong performance in isolation, yet practitioners lack guidance on which method works for their specific task. This pilot study compares semantic entropy and self-consistency under matched computational budgets---same model, same sample count, same benchmarks. The key finding is benchmark sensitivity: semantic entropy achieves AUROC 0.551 on HaluEval but 0.289 on TruthfulQA under identical conditions, while self-consistency performs near random on both (AUROC 0.444--0.474). These results suggest that method performance depends on benchmark-specific factors, and published evaluations on different benchmarks may not be directly comparable. The pilot sample size (N=20 per dataset) limits statistical power; confidence intervals are wide and overlapping. This work establishes a controlled comparison framework and identifies benchmark-method alignment as a consideration for hallucination detection evaluation. Full-scale validation is required before drawing definitive conclusions.
